\documentclass[12pt,a4paper]{article}

% =============================================================================
% Packages
% =============================================================================
\usepackage[utf8]{inputenc}
\usepackage[T1]{fontenc}
\usepackage[english]{babel}

% Math
\usepackage{amsmath,amssymb,amsthm}

% Graphics and figures
\usepackage{graphicx}
\usepackage{float}
\usepackage{subcaption}

% Tables
\usepackage{booktabs}
\usepackage{array}

% Typography
\usepackage{microtype}
\usepackage{lmodern}

% Layout
\usepackage{geometry}
\geometry{margin=1in}

% Links and references
\usepackage[colorlinks=true, linkcolor=blue, citecolor=blue, urlcolor=blue]{hyperref}
\usepackage{cleveref}

% Code listings (optional)
\usepackage{listings}
\lstset{
    basicstyle=\ttfamily\small,
    breaklines=true,
    frame=single,
    numbers=left,
    numberstyle=\tiny\color{gray},
}

% Colors
\usepackage{xcolor}

% Bibliography
\usepackage[style=numeric, backend=biber, sorting=nyt]{biblatex}
\addbibresource{references/references.bib}

% =============================================================================
% Theorem environments
% =============================================================================
\newtheorem{theorem}{Theorem}[section]
\newtheorem{lemma}[theorem]{Lemma}
\newtheorem{corollary}[theorem]{Corollary}
\newtheorem{definition}[theorem]{Definition}

% =============================================================================
% Document metadata
% =============================================================================
\title{Example Document: LaTeX Docker Environment}
\author{Your Name Here}
\date{\today}

% =============================================================================
% Document
% =============================================================================
\begin{document}

\maketitle

\begin{abstract}
    This is an example document to verify that the LaTeX Docker environment
    is working correctly. It demonstrates several common LaTeX features
    including math, figures, tables, code listings, and bibliography
    management with \texttt{biblatex} and \texttt{biber}.
\end{abstract}

\tableofcontents
\newpage

% --- Sections ---
\section{Introduction}
\label{sec:introduction}

Welcome to your LaTeX Docker environment! This document serves as both a
verification test and a reference for common LaTeX constructs.

\subsection{What This Environment Provides}

This setup includes:
\begin{itemize}
    \item \textbf{Full TeX Live distribution} — every package available
    \item \textbf{Multiple engines} — pdfLaTeX, XeLaTeX, LuaLaTeX
    \item \textbf{Automated builds} — via \texttt{latexmk} and Docker Compose
    \item \textbf{VS Code integration} — with LaTeX Workshop extension
    \item \textbf{Bibliography support} — using \texttt{biblatex} + \texttt{biber}
\end{itemize}

\subsection{Getting Started}

To compile this document, run:
\begin{lstlisting}[language=bash]
./compile.sh example
\end{lstlisting}

Or create a new project:
\begin{lstlisting}[language=bash]
./compile.sh new my-project
\end{lstlisting}

As shown in the literature~\cite{knuth1984texbook, lamport1994latex}, \LaTeX{}
remains the gold standard for scientific typesetting.

\section{Mathematics}
\label{sec:mathematics}

\LaTeX{} excels at typesetting mathematical content.

\subsection{Inline and Display Math}

Euler's identity $e^{i\pi} + 1 = 0$ is often considered the most beautiful
equation in mathematics.

The quadratic formula solves $ax^2 + bx + c = 0$:
\begin{equation}
    x = \frac{-b \pm \sqrt{b^2 - 4ac}}{2a}
    \label{eq:quadratic}
\end{equation}

We can reference \cref{eq:quadratic} easily with \texttt{cleveref}.

\subsection{Aligned Equations}

Maxwell's equations in differential form:
\begin{align}
    \nabla \cdot \mathbf{E}  &= \frac{\rho}{\varepsilon_0}
        \label{eq:gauss} \\
    \nabla \cdot \mathbf{B}  &= 0
        \label{eq:gauss-mag} \\
    \nabla \times \mathbf{E} &= -\frac{\partial \mathbf{B}}{\partial t}
        \label{eq:faraday} \\
    \nabla \times \mathbf{B} &= \mu_0 \mathbf{J}
        + \mu_0 \varepsilon_0 \frac{\partial \mathbf{E}}{\partial t}
        \label{eq:ampere}
\end{align}

\subsection{Theorem Environments}

\begin{theorem}[Pythagorean Theorem]
    \label{thm:pythagoras}
    For a right triangle with legs $a$, $b$ and hypotenuse $c$:
    \[ a^2 + b^2 = c^2 \]
\end{theorem}

\begin{proof}
    The proof is left as an exercise for the reader.
\end{proof}

\subsection{Matrices}

A rotation matrix in 2D:
\begin{equation}
    R(\theta) = \begin{pmatrix}
        \cos\theta & -\sin\theta \\
        \sin\theta &  \cos\theta
    \end{pmatrix}
\end{equation}

\section{Figures and Tables}
\label{sec:figures-tables}

\subsection{Tables}

\Cref{tab:engines} compares the \LaTeX{} engines available in this environment.

\begin{table}[H]
    \centering
    \caption{Comparison of \LaTeX{} engines}
    \label{tab:engines}
    \begin{tabular}{@{}llll@{}}
        \toprule
        \textbf{Engine} & \textbf{Font Support} & \textbf{Speed} & \textbf{Best For} \\
        \midrule
        pdfLaTeX  & Type1, OTF (limited) & Fast    & Most documents \\
        XeLaTeX   & OpenType, TrueType   & Medium  & Custom fonts, Unicode \\
        LuaLaTeX  & OpenType, TrueType   & Slower  & Lua scripting, fonts \\
        \bottomrule
    \end{tabular}
\end{table}

\subsection{Figures}

To include a figure, place your image in the \texttt{figures/} directory:

\begin{verbatim}
\begin{figure}[H]
    \centering
    \includegraphics[width=0.8\textwidth]{figures/my-image.pdf}
    \caption{Description of figure}
    \label{fig:example}
\end{figure}
\end{verbatim}

% Uncomment when you have an actual figure:
% \begin{figure}[H]
%     \centering
%     \includegraphics[width=0.8\textwidth]{figures/my-image.pdf}
%     \caption{An example figure}
%     \label{fig:example}
% \end{figure}

\section{Conclusion}
\label{sec:conclusion}

If this document compiled successfully, your \LaTeX{} Docker environment is
fully operational. You now have access to:

\begin{enumerate}
    \item The complete TeX Live distribution with all packages
    \item Three compilation engines: pdfLaTeX, XeLaTeX, and LuaLaTeX
    \item Automated compilation via \texttt{latexmk}
    \item Bibliography management with \texttt{biblatex} and \texttt{biber}
    \item VS Code integration with the LaTeX Workshop extension
\end{enumerate}

Happy typesetting!


% --- Bibliography ---
\newpage
\printbibliography

\end{document}
